\chapter{特征工程与数据预处理}

\section{学习目标}

学完本章后，你应该能够：

\begin{itemize}
    \item 解释为什么同一个模型会因为特征表达不同而表现差很多；
    \item 理解缺失值处理、标准化、类别编码和特征构造各自解决什么问题；
    \item 识别哪些特征更接近物理信息，哪些只是仪器或观测过程带来的“杂讯特征”；
    \item 明白预处理步骤必须只在训练集上拟合，避免数据泄漏；
    \item 理解为什么 \texttt{Pipeline} 和 \texttt{ColumnTransformer} 能帮助你把这些步骤组织得更可靠。
\end{itemize}

\section{为什么预处理会改变模型结果}

很多初学者会把机器学习理解成“换模型”。但在真实任务中，特征的组织方式往往和模型本身同样重要。即使模型完全不变，只要你改变：

\begin{itemize}
    \item 缺失值怎么填；
    \item 数值特征是否标准化；
    \item 类别变量怎么编码；
    \item 是否构造了更贴近物理问题的新特征；
\end{itemize}

模型的行为就可能发生明显变化。

\section{一个教学版预处理示例}

配套 notebook 使用 \nolinkurl{object_preprocessing_demo.csv}。这个数据集刻意包含了几类在真实项目中很常见的情况：

\begin{itemize}
    \item \textbf{缺失值}：部分样本缺少 \texttt{seeing\_arcsec} 或 \texttt{sky\_region}；
    \item \textbf{量纲尺度差异}：颜色指数大约在 $0$ 到 $2$ 之间，而 \texttt{detector\_x} 的量级在几百像素；
    \item \textbf{类别特征}：\texttt{sky\_region} 不是数值，而是一个离散标签；
    \item \textbf{可疑的杂讯特征}：\texttt{detector\_x} 更像仪器坐标，而不是物理上有意义的对象性质。
\end{itemize}

这正适合用来说明：预处理不是“美化数据”，而是在决定模型真正看到什么。

\section{缺失值：先承认，再处理}

缺失值不是异常情况，而是科学数据的常态。常见原因包括：

\begin{itemize}
    \item 某次观测失败；
    \item 某个波段没有被覆盖；
    \item 数据清洗时剔除了不可靠测量；
    \item 不同巡天或仪器之间字段不完整。
\end{itemize}

在教学例子里，我们对数值特征用训练集均值做最小填补，对类别特征用训练集众数填补。这不是唯一正确方法，但足够帮助学生建立一个关键原则：\textbf{填补参数必须由训练集决定，而不是偷看测试集。}

\section{标准化：让距离和权重变得可比较}

对于基于距离的模型，或者对特征尺度敏感的模型，不同量纲会直接改变模型的“几何形状”。如果一个无关特征的数值范围特别大，它可能会主导距离计算。

\begin{examplebox}[标准化的直觉]
\begin{lstlisting}[style=pythonstyle]
z = (x - mean_train) / std_train
\end{lstlisting}
\end{examplebox}

在 notebook 里，\texttt{detector\_x} 的量级远大于颜色指数。原始特征下，最近中心分类器会被这个量级巨大的列带偏；只要做了标准化，模型行为就会明显改善。

\section{类别编码：让模型读懂离散信息}

\texttt{sky\_region} 这种变量不能简单当作一个整数编号，因为：

\begin{itemize}
    \item \texttt{north} 不应该被解释成“小于” \texttt{south}；
    \item 编号之间的人为距离没有物理意义。
\end{itemize}

因此，最常见的起点是\textbf{one-hot 编码}：把每个类别拆成一个二元列。

\begin{examplebox}[one-hot 的最小形式]
\begin{lstlisting}[style=pythonstyle]
["equator", "north", "south"]
\rightarrow
[1, 0, 0]
\end{lstlisting}
\end{examplebox}

\section{特征工程：把物理直觉放进输入里}

特征工程不是随意堆砌更多列，而是尝试把领域知识转化成更有解释力的表示。在这个例子中，我们构造了一个简单的颜色和：

\begin{itemize}
    \item \texttt{color\_sum = u\_g + g\_r + r\_i}
\end{itemize}

它不一定是最佳特征，但足以让学生看到：有时一个经过整理的组合特征，比机械地把原始列全部丢进模型更清楚。

\section{为什么有些特征应该被丢掉}

\texttt{detector\_x} 是本章中一个故意设计出来的教学“陷阱”。它在数值上很大，容易影响距离，但从科学意义上看，它更接近观测器坐标，而不是对象的天体物理属性。

这提醒我们：不是所有“可用列”都值得保留。特征选择并不只是统计操作，也是在问：

\begin{itemize}
    \item 这个特征是否真的携带物理信息？
    \item 它会不会只是仪器、管线或观测条件的痕迹？
    \item 它会不会导致模型学到错误但表面高分的规律？
\end{itemize}

\section{Pipeline 与 ColumnTransformer 为什么重要}

当预处理步骤变多时，手工操作很容易出错。例如：

\begin{itemize}
    \item 先填补、后标准化，还是先标准化、后填补？
    \item 哪些列做 one-hot，哪些列做数值缩放？
    \item 测试集是不是误用了全数据的均值？
\end{itemize}

这正是 \texttt{Pipeline} 和 \texttt{ColumnTransformer} 的价值所在：它们能把预处理和模型训练明确串联起来，减少手工重复与数据泄漏风险。

\section{最需要警惕的不是技术错误，而是数据泄漏}

预处理里最隐蔽的风险之一，就是在划分训练集和测试集之前就先对全数据做：

\begin{itemize}
    \item 均值填补；
    \item 标准化；
    \item 特征选择；
    \item 编码规则拟合。
\end{itemize}

这样一来，测试集的信息就提前泄漏给了模型。哪怕模型结构本身没变，评估结果也会被污染。

\begin{warnbox}[预处理中的常见误区]
\begin{itemize}
    \item 用全数据的均值填补缺失值；
    \item 忽视量纲差异，让大尺度噪声特征主导模型；
    \item 把类别变量直接编码成有序整数；
    \item 机械增加特征数量，而不判断它们是否有物理意义；
    \item 在没有清楚记录步骤的情况下手工反复处理数据。
\end{itemize}
\end{warnbox}

\section{AI 助手如何帮助你}

在这一章，AI 助手很适合帮助你：

\begin{itemize}
    \item 检查你的缺失值处理和标准化代码是否只使用了训练集统计量；
    \item 讨论某个特征是否更像物理信号还是仪器痕迹；
    \item 把手写预处理流程整理成更清楚的函数或 \texttt{Pipeline} 结构。
\end{itemize}

但哪些特征真正值得保留、是否存在隐蔽泄漏、某个工程特征是否合理，仍然必须由你结合科学背景来判断。

\section{本章小结}

特征工程与数据预处理不是模型外面的“准备工作”，而是模型定义本身的一部分。你让模型看到什么、用什么尺度看到、缺失值怎么补、类别怎样编码，都会直接决定它学到的规律是否可信。
